{\tiny\setlength{\tabcolsep}{3pt}
\begin{longtable}{p{2.1cm} p{3cm} p{1.3cm} p{3cm} p{2.9cm} p{3.4cm}}
\caption{Mandate landscape, 24 jurisdictions. URLs and quotes are in the source extraction table supplied with the data.}\label{tab:mandates}\\
\toprule
\textbf{Jurisdiction} & \textbf{Instrument} & \textbf{Year} & \textbf{Applies to} & \textbf{Penalty / incentive} & \textbf{Reported compliance}\\ \midrule\endfirsthead
\toprule \textbf{Jurisdiction} & \textbf{Instrument} & \textbf{Year} & \textbf{Applies to} & \textbf{Penalty / incentive} & \textbf{Reported compliance}\\ \midrule\endhead
\bottomrule\endfoot
India (Model Building Bye-Laws 2016, MoHUA) & Model Building Bye-Laws 2016 (MBBL), MoHUA; rainwater harvesting chapter & 2016 & all buildings having a plot size of 100 sq.m or more (RWH proposal required with building plan) & NR (states decide) & 33 States/UTs have adopted the RWH provisions (Lok Sabha reply 05.03.2020; also PIB 24.03.2020); 'barring Sikkim, Lakshadweep and Mizoram' all States/UTs adopted (Lok Sabha Q1241, 09.02.2023)\\
Tamil Nadu (incl. Chennai) & Tamil Nadu Municipal Laws (Second Amendment) Ordinance, 2003 (TN Ordinance No. 4 of 2003) inserting s.255-A Chennai City Municipal Corporation Act etc. & 2003 (promulgated 19 July 2003) & every owner or occupier of a building (all buildings; Govt buildings too) & water supply connection disconnected until structure provided; Commissioner may build and recover cost as property tax; new buildings not given water/sewer connection without RWH (CSE) & Chennai 2019: GCC inspected \textasciitilde 2.72 lakh houses since Aug 2019, \textasciitilde 1.62 lakh had RWH (approx 60\%), 9,000 notices issued; CMWSSB claims 8.9 lakh structures; Rain Centre 2015 audit of 1,100 samples found govt buildings worst violators\\
Bengaluru (BWSSB jurisdiction) & Bangalore Water Supply and Sewerage (Amendment) Act, 2011 (Karnataka Act No. 05 of 2011) substituting s.72A (bill introduced 2009); penalty regulations 2015; BWSS (Amendment) Bill 2021 extends to 30x40 sites & 2009 (bill) / 2011 (Act assented 2 Feb 2011) / 2021 (extension) & existing buildings on sital area >= 2400 sq ft (60x40 ft site); new buildings on sital area >= 1200 sq ft (30x40 ft); 2021 Bill: sital area 2,325-10,763.9 sq ft RWH; >=10,763.9 sq ft dual piping + RWH & residential: 25\% of water bill first 3 months then 50\%; non-residential: 50\% then 100\% (Citizen Matters 2023); Rs 1.95-2 crore/month collected & Dec 2023: \textasciitilde 1.9 lakh of 10.8 lakh properties with water connections (\textasciitilde 18\%) have RWH; 39,703 properties paying monthly penalty; Mar 2024 (The Hindu): 1,96,591 residential buildings implemented, 39,213 have not\\
Karnataka (statewide, outside BWSSB) & Government orders for Government buildings; statewide urban mandate 'contemplated' (2011); 20\% tax rebate (CSE) & NR & Government buildings; amendment of bye-laws in cities >20 lakh population under action (CSE); statewide new buildings contemplated 2011 & 20\% rebate on tax payment for 5 years (CSE) & NR\\
Delhi (NCT) & Building Bye-laws 1983 modified by MoUD Gazette Notification 28.07.2001 (plots >=100 sq m); CGWA directive; DJB resolution No. 829; Unified Building Bye-Laws 2016 & 2001 (bye-law), 2012 (DJB mandatory per Outlook), 2016 (UBBL) & all buildings with plot area >= 100 sq m (2001 notification; DJB resolution 829); DJB compulsory for plots >= 500 sq m; UBBL 2016 features mandatory for new constructions above 105 sq m (DUAC) & building plans not sanctioned without RWH; DJB 10\% rebate on water bill for plots >= 100 sq m with functional RWH; penalty 1.5 times water bill / '50\% of the monthly bill amount' for non-compliance; proposed environmental compensation Rs 50,000 to Rs 5 lakh (2023 proposal); NGT Rs 5 lakh per educational institution (OA 217/2016) & Delhi: 1,560 consumers availed DJB rebate (NGT monitoring committee report OA 496/2016); 3,675 govt schools/colleges of 4,769 institutions had functional systems; only 1,869 private buildings as of 2021 (Citizen Matters); Dwarka: NGT expert committee inspected 354 CGHS, 'Most structures are found either completely defunct or just partially working' (2023)\\
Mumbai / Maharashtra (Greater Mumbai DCR; all towns non-gaothan) & Govt of Maharashtra UDD Notification No. TPB-4307/396/CR-124/2007/UD-11 dated 6 June 2007 modifying Development Control Regulations for Greater Mumbai 1991; BMC mandatory for new buildings since 2002 (Citizen Matters); Aug 2026 GR: old buildings seeking new water/drainage connections must have RWH & 2002 (BMC, new buildings) / 2007 (300 sq m notification) & plots having area not less than 300 sq m in non-gaothan areas of all towns (2007); earlier 1,000 sq m (CSE) & no building permission unless provision for RWH; levy not exceeding Rs 1000 per annum per 100 sq m built-up area for failure to provide/maintain; completion/occupation certificate only after RWH (CSE) & approx. 3,000 societies of an estimated 20,000 had RWH 2007-2021 (Citizen Matters Jul 2024)\\
Hyderabad / Telangana (HMWSSB area) & AP Govt order (2000) for all municipal corporations/urban areas (IWP); APWALTA 2002; HMWSSB directive dated 14 Dec 2024 for 17,000 consumers & 2000 (GO; tentative enforcement June 2001) & all new buildings with area 300 sq m or more irrespective of roof area; existing buildings within one year of GO (CSE); originally 200 sq m then 300 sq m (Realty Today) & 50\% subsidy on RWH structures was lifted due to poor response (CSE); 2024: double charges for HMWSSB water tankers for non-compliant (Rs 500 to Rs 1,000 per 5 kL) & Dec 2024: NGO survey of 39,000 properties; 42,000 buildings identified needing RWH; 22,000 installed; 17,000 non-compliant (Realty Today / South First)\\
Pune (PMC) & PMC requirement for new buildings (2006, per Realty Today); property-tax rebate scheme & 2006 & all new buildings (Realty Today) & property tax rebate up to 10\% for RWH (PMC eco-housing rebate) & 306 properties availed rebate in FY2012-13; 3,716 in FY2022-23, of approx. 12.5 lakh taxable properties (approx 0.3\%)\\
Nagpur (NMC) & NMC regulation (mandatory from March 2005 per CSE; ruling effective 28 Aug 2009 per Propnewstime) & 2005 (CSE) / 2009 (Propnewstime) & new constructions of area >= 300 sq m (non-gaothan); housing society layouts and amenity spaces & fine up to Rs 1,000 per annum per 100 sq m built-up area; NMC 5-10\% property tax discount for eco initiatives & May 2023: only 1,123 buildings comply = 0.18\% of \textasciitilde 6.36 lakh properties\\
Gujarat / Ahmedabad (AUDA/AMC GDCR) & Gujarat Development Control Regulations (GDCR); AUDA notification 2002 (1,500 sq m+); State Govt buildings after 2003 & 2002 (AUDA) / 2003 (Govt buildings) & GDCR: buildings 500-1,500 sq m follow RWH specifications; 1,500-4,000 sq m percolation wells (one per 4,000 sq m); AUDA 2002: constructions >= 1,500 sq m & Surat: 50\% subsidy up to Rs 2,000 (CSE); no penalty stated & NR ('still not very popular', 2008 clipping)\\
Rajasthan / Jaipur (JDA) & Rajasthan state mandate (PHED certification); Jaipur/Rajasthan Building Bye-Laws 2020 Regulation 21.7 (rule 10.11.1); JDA enforcement 2011 & 2011 (JDA penalty) / 2020 (Bye-laws) & all buildings (any land use) with plot area over approx. 300 sq m (CSE legislation page) / over 500 sq m (CSE laws page); Bye-laws 2020: each plot owner (no minimum stated in translation) & water supply connection can be terminated if RWH completion certificate not submitted to PHED; JDA penalty Rs 25,000 (2011); Bye-laws 2020: annual maintenance fee levied by local authority & Jul 2023: only 220 buildings in Jaipur had complied (JDA)\\
Madhya Pradesh / Indore (also Jabalpur, Gwalior) & MP Government Order no. F 23 (107) 95 XXXII dated 7 July 2004; MP Bhumi Vikas Rules 1984; Jabalpur from 1 Apr 2005; Gwalior order 27 Jan 2006 & 2004 & all new buildings with area >= 250 sq m; mandatory for G+3 structures & 6\% property tax rebate; Gwalior penalty Rs 7,000 for non-compliance & NR\\
Kerala & Kerala Municipality Building (Amendment) Rules, 2004 (Extraordinary Gazette No. 92/2004 dated 12-01-2004) inserting Chapter XVI-A, Rule 109A & 2004 & all new buildings: Residential A1 with floor area >= 100 sq m and plot area >= 200 sq m; A2, B, C, D, E and specified G1/G2 occupancies & no penalty stated; exemption where water logging common; KWA 90\% / Jalanidhi 75\% subsidies (CSE) & NR (CSE: flats with multiple owners maintained better than individual houses)\\
Haryana (incl. Gurugram, Faridabad) & HUDA rules (CSE); Haryana Building Code 2017 & 2017 (Code) & roof rainwater harvesting where rooftop area >= 100 sq m; recharge for Group Housing with plot > 500 sq m (Tribune); HUDA: all new buildings irrespective of roof area (CSE) & NR & Faridabad: MCF survey 2018-19 found \textasciitilde 75-80\% of RWH systems defunct (Tribune Dec 2023)\\
Uttar Pradesh (Kanpur; Ghaziabad; statewide CM direction 2025) & Building bye-laws (Kanpur: >=1000 sq m per CSE legislation page / >=500 sq m per CSE laws page); Ghaziabad bye-laws >=300 sq m; CM direction 4 Oct 2025 for buildings >100 sq m & NR (bye-laws); 2025 (CM order) & Kanpur: new buildings >= 1000 sq m (or 500 sq m); Ghaziabad: plots >= 300 sq m; CM 2025: buildings larger than 100 sq m, urban and rural & Ghaziabad: penalties up to Rs 3 lakh depending on plot size and severity (2025) & Ghaziabad: of 100 identified locations, 20 visited (Jul 2025); no counts\\
Punjab (incl. Ludhiana) & Punjab building bylaws notified 2018 (Tribune): RWH for structures > 119 sq yards (earlier 250 sq yd residential / 500 sq yd commercial); RWH Cell in each ULB & 2018 & structures (residential, commercial, institutional) on more than 119 sq yards (approx 100 sq m); earlier residential > 250 sq yd and commercial > 500 sq yd & utility connections denied without RWH and completion certificate (Tribune Bathinda) & NR ('only several property owners ensured proper installation and functioning')\\
Chandigarh (UT) & Chandigarh Administration order (2008) for plots of one kanal or more; CSE: all new buildings & 2008 & plots of one kanal (approx 500 sq yd / 418 sq m) or more; administration mulling lower threshold (2019) & NR (no penalty stated; 'absence of any mechanism to keep a check') & 'only a handful of private building owners complying with it' (Pioneer 2019)\\
West Bengal / Kolkata & West Bengal Municipal (Building) Rules, 2007, Rule 171 & 2007 & NR (threshold not stated in source) & NR & NR\\
Goa & Policy on Rainwater Harvesting, Govt of Goa (15/10/2008, WRD); RWH policy notified in gazette Dec 2022 with plot thresholds (Propnewstime) & 2008 (policy) / 2022 (gazette policy) & 2022 policy: residential complexes/apartments on plots >= 2,000 sq m; commercial complexes >= 1,500 sq m; industrial units >= 10,000 sq m; all government buildings; exempt if STP/greywater recycling exists; 2008: PWD govt buildings (CSE) & NR & NR\\
Himachal Pradesh & State direction: no building plan approved in urban areas without RWH (CSE) & NR & all buildings to be constructed in urban areas; schools, government buildings, rest houses, industries, bus stands & building plan not approved without RWH system & NR\\
Uttarakhand (Mussoorie-Dehradun Development Authority) & MDDA building byelaws provision (CSE); MDDA mandatory for new houses of 1,200 sq ft (Tribune 2011) & NR & all new houses of 1,200 sq ft (Tribune, Mussoorie) & NR & Mussoorie 2011: 'only 30\% hoteliers have adopted roof-top rain water harvesting'\\
Andhra Pradesh (all municipal corporations/urban areas; Visakhapatnam) & AP Government Order (2000) making RWH mandatory in all municipal corporations, UDAs and municipalities; APWALTA 2002 & 2000 & buildings of 300 sq m or more & NR & NR (no Visakhapatnam-specific data found)\\
Rajkot (Gujarat) & Rajkot General Development Control Regulations, Regulation No. 17.18 & NR & buildings with built-up area >= 80 sq m; existing buildings > 1,000 sq m five-year moratorium & NR & NR\\
Ranchi (Jharkhand) & Ranchi Regional Development Authority byelaws & NR & NR & NR & NR\\
\end{longtable}}
